\documentclass[11pt,letterpaper]{article}
\usepackage[spanish]{babel}
\usepackage[utf8]{inputenc}
\usepackage[T1]{fontenc}
\usepackage{geometry}
\usepackage{amsmath}
\usepackage{booktabs}
\usepackage{hyperref}
\usepackage{enumitem}

\geometry{margin=1in}
\setlist[itemize]{leftmargin=1.2em}

\title{RiskGate: Nueva Versión del Modelo de Opciones}
\author{RiskGate}
\date{Mayo 2026}

\begin{document}
\maketitle

\section*{Resumen}

Esta versión mantiene la filosofía original de RiskGate: operación local, datos de mercado de solo lectura, sin scraping de broker, sin colocación de órdenes y sin trading automático. El objetivo sigue siendo apoyar la disciplina antes de entrar a una operación, no predecir precios ni dar asesoría financiera.

La actualización hace que el modelo sea más correcto para opciones. Ahora puede evaluar estructuras multi-leg, calcular indicadores técnicos desde barras OHLCV, medir la calidad de contratos con funciones continuas, simular escenarios de payoff con Monte Carlo y resumir narrativa pública de Reddit como una señal cualitativa opcional.

\section{Estructura Multi-Leg}

Antes, una idea de trade se aproximaba con un solo contrato principal. Eso funcionaba para una call o put simple, pero era insuficiente para spreads, coberturas, collars o estrategias con piernas largas y cortas.

La nueva tabla \texttt{OptionLeg} guarda cada pierna con:

\begin{itemize}
  \item tipo de opción: call o put;
  \item lado: long o short;
  \item cantidad, vencimiento y strike;
  \item bid, ask, mid y spread;
  \item IV, IV rank, Greeks, volumen, open interest y DTE.
\end{itemize}

RiskGate conserva los campos antiguos para compatibilidad. Si una idea vieja no tiene piernas explícitas, el sistema genera una pierna virtual a partir del contrato legacy.

\section{Métricas de Estrategia}

El nuevo motor de payoff calcula la utilidad o pérdida al vencimiento por pierna:

\[
\text{callPayoff} = \max(S_T - K, 0)
\]

\[
\text{putPayoff} = \max(K - S_T, 0)
\]

Para una pierna long:

\[
\text{PnL} = q \times 100 \times (\text{payoff} - \text{prima pagada})
\]

Para una pierna short:

\[
\text{PnL} = q \times 100 \times (\text{prima recibida} - \text{payoff})
\]

La suma de las piernas produce métricas de estrategia:

\begin{itemize}
  \item débito neto y crédito neto;
  \item pérdida máxima;
  \item ganancia máxima;
  \item puntos de break-even;
  \item Delta, Gamma, Theta y Vega totales;
  \item curva determinística de payoff.
\end{itemize}

El modelo bloquea por defecto piernas short desnudas cuando no existe una cobertura definida por otra pierna long compatible.

\section{Indicadores Técnicos Automáticos}

La nueva tabla \texttt{PriceBar} almacena barras históricas por portafolio, símbolo, temporalidad y timestamp. Con esas barras se calculan:

\begin{itemize}
  \item SMA20, SMA50 y SMA200;
  \item RSI14;
  \item ATR14;
  \item volumen relativo;
  \item BWAP;
  \item distancia del precio al BWAP.
\end{itemize}

BWAP se implementa como un promedio ponderado por volumen estilo VWAP:

\[
\text{BWAP} =
\frac{\sum \text{precio típico} \times \text{volumen}}
{\sum \text{volumen}}
\]

donde:

\[
\text{precio típico} = \frac{high + low + close}{3}
\]

Si no existen high y low, el cálculo usa close por volumen.

\section{Score Técnico Continuo}

El score técnico ya no depende solo de etiquetas manuales como ``arriba de la media de 20 días''. Ahora combina calidades continuas:

\[
\text{technicalScore} = 100 \times (
0.25Q_{trend} +
0.20Q_{MA} +
0.15Q_{RSI} +
0.15Q_{volume} +
0.10Q_{BWAP} +
0.10Q_{SR} +
0.05Q_{ATR})
\]

Para trades bullish, favorece precio sobre SMA20/SMA50, SMA20 sobre SMA50, cierre sobre BWAP, RSI saludable y volumen relativo fuerte. Para trades bearish, invierte la lógica de medias y BWAP. Para trades neutrales, favorece precio cerca de BWAP y menor expansión de ATR.

El usuario puede mantener inputs manuales, pero RiskGate muestra advertencias cuando esos inputs contradicen los indicadores calculados.

\section{Calidad Continua de Opciones}

La calidad de opciones usa una función continua con desglose por componente:

\[
\text{optionsQuality} = 100 \times (
0.30Q_{spread} +
0.20Q_{DTE} +
0.15Q_{volume} +
0.15Q_{OI} +
0.08Q_{IVRank} +
0.07Q_{premium} +
0.05Q_{theta})
\]

\begin{table}[h]
\centering
\begin{tabular}{ll}
\toprule
Componente & Lógica \\
\midrule
Spread & Crédito completo hasta 2\%; decae suavemente después. \\
DTE & Cero bajo 5 días; máximo entre 21 y 60 días. \\
Volumen & Escala con raíz cuadrada hacia 500 contratos. \\
Open interest & Escala con raíz cuadrada hacia 1000 contratos. \\
IV rank & Penalización cuadrática para IV rank alto. \\
Prima & Máxima calidad bajo 0.25\% del portafolio. \\
Theta & Penaliza theta alto frente al mid del contrato. \\
\bottomrule
\end{tabular}
\end{table}

Los hard stops existentes se mantienen: spread mayor a 15\%, DTE menor a 7, prima mayor a 1\% del portafolio y liquidez muy baja.

\section{Monte Carlo}

El motor Monte Carlo simula precios terminales con movimiento browniano geométrico:

\[
S_T = S_0 \exp((\mu - q - 0.5\sigma^2)\tau + \sigma\sqrt{\tau}Z)
\]

Por defecto:

\begin{itemize}
  \item drift esperado \(\mu = 0\);
  \item tasa libre de riesgo \(r = 0.04\);
  \item dividend yield \(q = 0\);
  \item 10,000 simulaciones en el módulo base.
\end{itemize}

Cada simulación calcula el payoff de toda la estructura multi-leg. El resultado muestra probabilidad de utilidad, valor esperado, mediana, P5, P95, CVaR95, pérdida máxima simulada, ganancia máxima simulada e histograma.

Monte Carlo no reemplaza los hard stops. Solo agrega un overlay pequeño:

\begin{itemize}
  \item +5 si EV es positivo y la probabilidad de utilidad supera 55\%;
  \item -5 si EV es negativo;
  \item -10 si CVaR95 está cerca de la pérdida máxima.
\end{itemize}

\section{Reddit RAG y Sentimiento}

El módulo de Reddit está apagado por defecto. Cuando se configuran credenciales, RiskGate usa datos públicos de Reddit, limpia el texto y guarda solo campos necesarios para auditoría local.

El análisis clasifica la narrativa como:

\begin{itemize}
  \item strong bullish;
  \item bullish;
  \item mixed;
  \item bearish;
  \item strong bearish;
  \item hype;
  \item fear;
  \item low signal.
\end{itemize}

La señal de Reddit no aprueba trades. Puede sumar poco si coincide con la dirección y tiene confianza alta, pero penaliza hype, miedo extremo o preocupaciones recurrentes como fraude, dilución, lawsuit, riesgo de earnings, colapso de demanda o valuación extrema.

\section{Arquitectura del Score}

El score base conserva los pesos adaptativos:

\[
\text{baseScore} =
M w_M + T w_T + O w_O + Q w_Q
\]

Después se agregan overlays:

\[
\text{finalScore} =
\text{clamp}(\text{baseScore} + \text{MonteCarloOverlay} + \text{SentimentOverlay}, 0, 100)
\]

Los hard stops siguen anulando todo:

\[
\text{si existe hard stop} \Rightarrow \text{decisión} = reject
\]

\section{Sizing de Riesgo}

La aprobación no implica tamaño completo. La nueva fórmula de sizing es:

\[
\text{risk} =
\text{baseRisk}
\times C
\times L
\times R
\times V
\times K
\times MC
\times S
\]

donde:

\begin{itemize}
  \item \(C\): multiplicador de confianza;
  \item \(L\): multiplicador de liquidez;
  \item \(R\): multiplicador de régimen de mercado;
  \item \(V\): multiplicador de volatilidad;
  \item \(K\): multiplicador de concentración;
  \item \(MC\): multiplicador de Monte Carlo;
  \item \(S\): multiplicador de sentimiento.
\end{itemize}

Luego se aplican los límites finales del portafolio:

\[
\text{risk final} = \min(
\text{risk},
\text{sleeve restante},
\text{límite mensual restante},
\text{riesgo abierto restante})
\]

\section{Conclusión}

La nueva versión convierte a RiskGate en un sistema más robusto para decidir y dimensionar trades de opciones. La mejora principal no es ``predecir mejor'', sino evitar errores estructurales: spreads mal medidos, contratos ilíquidos, riesgo short indefinido, señales técnicas subjetivas y tamaños de posición demasiado agresivos.

RiskGate sigue siendo una herramienta de apoyo a decisiones. La ejecución queda fuera del sistema y requiere juicio humano.

\end{document}
